\documentclass[10pt,a4paper]{article}
\usepackage[a4paper,margin=14mm,top=12mm,bottom=14mm]{geometry}
\usepackage{fontspec}
\setmainfont{Liberation Sans}
\usepackage{polyglossia}
\setdefaultlanguage{french}
\usepackage{amsmath}
\usepackage{xcolor}
\usepackage{graphicx}
\usepackage[most]{tcolorbox}
\usepackage{tabularx,booktabs,array,multirow}
\usepackage{enumitem}
\usepackage{fancyhdr}
\usepackage{lastpage}
\usepackage[hidelinks]{hyperref}
\definecolor{uimm}{HTML}{005677}
\definecolor{cyan}{HTML}{13A2DD}
\definecolor{fond}{HTML}{E6F2FA}
\definecolor{ok}{HTML}{2E7D32}
\definecolor{ko}{HTML}{B3261E}
\definecolor{violet}{HTML}{7B1FA2}
\setlength{\parindent}{0pt}
\setlength{\parskip}{3pt}
\setlist{nosep}
\setlist[itemize]{leftmargin=13pt,label=\textcolor{cyan}{\textbullet}}
\newcommand{\titre}[1]{\vspace{5pt}{\color{uimm}\large\bfseries #1}\par\vspace{-3pt}{\color{cyan}\rule{\linewidth}{1.2pt}}\par\vspace{0pt}}
\tcbset{boite/.style={colback=fond,colframe=uimm,boxrule=0.6pt,arc=2pt,left=6pt,right=6pt,top=4pt,bottom=4pt},
        defi/.style={colback=violet!6,colframe=violet,boxrule=0.8pt,arc=2pt,left=6pt,right=6pt,top=4pt,bottom=4pt,
                     title={\bfseries Défi},coltitle=white,colbacktitle=violet,fonttitle=\small}}
\newcommand{\bandeau}[3]{%
\begin{tcolorbox}[colback=uimm,colframe=uimm,arc=2pt,left=4pt,right=8pt,top=4pt,bottom=4pt]
\begin{minipage}[c]{18mm}\setlength{\fboxsep}{1.2pt}\fcolorbox{white}{white}{\includegraphics[height=16.5mm]{logo_uimm.jpg}}\end{minipage}\hspace{4mm}%
\begin{minipage}[c]{\dimexpr\linewidth-22mm\relax}\color{white}
{\small #1}\\[1pt]{\LARGE\bfseries #2}\\[2pt]{\large #3}\hfill{\small CQPM Technicien qualité}
\end{minipage}
\end{tcolorbox}}
\newcommand{\reponse}[1][1]{\par\vspace{2pt}{\color{gray!60}\foreach \i in {1,...,#1}{\rule{\linewidth}{0.3pt}\par\vspace{9pt}}}}
\usepackage{pgffor}
\pagestyle{fancy}\fancyhf{}
\renewcommand{\headrulewidth}{0pt}
\newcommand{\pied}[1]{\fancyfoot[L]{\footnotesize\color{uimm}#1}\fancyfoot[R]{\footnotesize\color{uimm}Page \thepage{} sur \pageref{LastPage}}}

\usepackage{tikz}
\usepackage{longtable}
\renewcommand{\bandeau}[3]{%
\begin{tcolorbox}[colback=uimm,colframe=uimm,arc=2pt,left=4pt,right=8pt,top=4pt,bottom=4pt]
\begin{minipage}[c]{18mm}\setlength{\fboxsep}{1.2pt}\fcolorbox{white}{white}{\includegraphics[height=16.5mm]{logo_uimm.jpg}}\end{minipage}\hspace{4mm}%
\begin{minipage}[c]{\dimexpr\linewidth-22mm\relax}\color{white}
{\small #1}\\[1pt]{\LARGE\bfseries #2}\\[2pt]{\large #3}
\end{minipage}
\end{tcolorbox}}
% en-tête d'une fiche : code, titre, question, bloc, durée, organisation, niveau
\newcommand{\entete}[7]{%
\begin{tcolorbox}[colback=uimm,colframe=uimm,arc=2pt,left=4pt,right=6pt,top=4pt,bottom=4pt]
\begin{minipage}[c]{18mm}\setlength{\fboxsep}{1.2pt}\fcolorbox{white}{white}{\includegraphics[height=16.5mm]{logo_uimm.jpg}}\end{minipage}\hspace{3mm}%
\begin{minipage}[c]{\dimexpr\linewidth-58mm\relax}\color{white}\raggedright
{\small Base SPC · #4}\\[1pt]{\Large\bfseries #1 · #2}\\[2pt]{\itshape #3}
\end{minipage}\hfill
\begin{minipage}[c]{33mm}\color{white}\small\raggedleft
Durée : \textbf{#5}\\ \textbf{#6}\\ Niveau #7 sur 3
\end{minipage}
\end{tcolorbox}}
\newcommand{\tagE}{{\setlength{\fboxsep}{1.5pt}\colorbox{cyan}{\color{white}\scriptsize\bfseries Ellistat}}}
\newcommand{\tagW}{{\setlength{\fboxsep}{1.5pt}\colorbox{uimm}{\color{white}\scriptsize\bfseries Page web}}}
\tcbset{pred/.style={colback=orange!6,colframe=orange!80!black,boxrule=0.6pt,arc=2pt,left=6pt,right=6pt,top=3pt,bottom=1pt,
                    title={\bfseries Je prédis, avant de commencer},coltitle=white,colbacktitle=orange!80!black,fonttitle=\small},
        retenir/.style={colback=fond,colframe=uimm,boxrule=0.6pt,arc=2pt,left=6pt,right=6pt,top=3pt,bottom=3pt,
                    title={\bfseries À retenir},coltitle=white,colbacktitle=uimm,fonttitle=\small},
        plus/.style={colback=violet!6,colframe=violet,boxrule=0.6pt,arc=2pt,left=6pt,right=6pt,top=3pt,bottom=3pt,
                    title={\bfseries Pour aller plus loin},coltitle=white,colbacktitle=violet,fonttitle=\small},
        act/.style={colback=uimm!10,colframe=uimm,boxrule=0.5pt,arc=2pt,left=5pt,right=5pt,top=2pt,bottom=2pt,before skip=7pt,after skip=3pt},
        verif/.style={colback=orange!6,colframe=orange!80!black,boxrule=0.6pt,arc=2pt,left=6pt,right=6pt,top=3pt,bottom=3pt,
                    title={\bfseries À vérifier sur poste avant la première séance},coltitle=white,colbacktitle=orange!80!black,fonttitle=\small}}
\newcommand{\fichier}[1]{\texttt{\small #1}}
\newcommand{\menu}[1]{\textsf{\textbf{#1}}}
\newcommand{\champ}[1]{\rule{#1}{0.3pt}}

\pied{Base SPC · fiche M4 · Pied à coulisse contre micromètre}
\begin{document}
\thispagestyle{fancy}
\entete{M4}{Pied à coulisse contre micromètre}{Deux instruments, une cale : lequel croire ?}{Le moyen de mesure}{30 min}{en binôme}{2}
\begin{tcolorbox}[boite]\textbf{Le but.} Comparer deux instruments sur la même cale : résolution, dispersion, décalage.\par\textbf{Il vous faut.} Un micromètre 0-25 mm · Un pied à coulisse au 1/50 · Une cale étalon de 10 mm\end{tcolorbox}
\begin{tcolorbox}[pred]Lequel des deux instruments donnera les mesures les plus groupées ? De combien, à votre avis ?\reponse\end{tcolorbox}
\titre{À faire}
\begin{enumerate}[label=\textcolor{cyan}{\textbf{\arabic*.}},leftmargin=15pt,itemsep=1.5pt]
\item Mesurer la cale 15 fois au micromètre (ou reprendre les mesures de M3), puis 15 fois au pied à coulisse. Retirer la cale à chaque fois.
\item \tagE\ \menu{Gage R\&R}, onglet \menu{Justesse}, « Ajouter une justesse » : une étude par instrument. Nombre de mesures 15 ; Valeur vraie = valeur de la cale ; LSI = cale $-$ 0,075 et LSS = cale + 0,075 ; k 3. Saisir toutes les décimales lues : 10,00 et 10,02 au pied à coulisse, 10,000 au micromètre.
\item \tagW\ Coller les deux séries dans la page « Comparer deux séries de mesures ».
\end{enumerate}
\vspace{3pt}\begin{minipage}[t]{0.485\linewidth}\textbf{Micromètre (mm)}\par\vspace{2pt}
{\renewcommand{\arraystretch}{1.4}\begin{tabularx}{\linewidth}{|*{3}{>{\centering\arraybackslash}p{5mm}|>{\centering\arraybackslash}X|}}\hline
\scriptsize 1 &  & \scriptsize 2 &  & \scriptsize 3 & \\\hline
\scriptsize 4 &  & \scriptsize 5 &  & \scriptsize 6 & \\\hline
\scriptsize 7 &  & \scriptsize 8 &  & \scriptsize 9 & \\\hline
\scriptsize 10 &  & \scriptsize 11 &  & \scriptsize 12 & \\\hline
\scriptsize 13 &  & \scriptsize 14 &  & \scriptsize 15 & \\\hline
\end{tabularx}}\end{minipage}\hfill\begin{minipage}[t]{0.485\linewidth}\textbf{Pied à coulisse (mm)}\par\vspace{2pt}
{\renewcommand{\arraystretch}{1.4}\begin{tabularx}{\linewidth}{|*{3}{>{\centering\arraybackslash}p{5mm}|>{\centering\arraybackslash}X|}}\hline
\scriptsize 1 &  & \scriptsize 2 &  & \scriptsize 3 & \\\hline
\scriptsize 4 &  & \scriptsize 5 &  & \scriptsize 6 & \\\hline
\scriptsize 7 &  & \scriptsize 8 &  & \scriptsize 9 & \\\hline
\scriptsize 10 &  & \scriptsize 11 &  & \scriptsize 12 & \\\hline
\scriptsize 13 &  & \scriptsize 14 &  & \scriptsize 15 & \\\hline
\end{tabularx}}\end{minipage}\par
\titre{J'observe, je décide}
\begin{enumerate}[label=\textcolor{cyan}{\textbf{Q\arabic*.}},leftmargin=22pt]
\item Micromètre : Cg = \champ{12mm} Cgk = \champ{12mm} \quad Pied à coulisse : Cg = \champ{12mm} Cgk = \champ{12mm}\par\vspace{5pt}
\item Quelle ligne des résultats d'Ellistat met en cause le pied à coulisse avant même de regarder Cg ?\reponse[1]
\item L'écart-type du pied à coulisse vaut combien de fois celui du micromètre ?\reponse[1]
\item Le pied à coulisse est-il un mauvais instrument ? Pour quelle tolérance conviendrait-il ?\reponse[2]
\end{enumerate}
\begin{tcolorbox}[retenir]Un instrument n'est ni bon ni mauvais : il est capable ou non \textbf{pour une tolérance donnée}.\end{tcolorbox}
\vfill{\footnotesize\color{uimm}\textbf{Ressources de la base :} Comparer deux séries de mesures · Données de secours.}
\end{document}